%-------------------------
% Resume in Latex
% Author : Jake Gutierrez
% Based off of: https://github.com/sb2nov/resume
% License : MIT
%------------------------

\documentclass[letterpaper,11pt]{article}

\usepackage{latexsym}
\usepackage[empty]{fullpage}
\usepackage{titlesec}
\usepackage{marvosym}
\usepackage[usenames,dvipsnames]{color}
\usepackage{verbatim}
\usepackage{enumitem}
\usepackage[hidelinks]{hyperref}
\usepackage{contour}
\usepackage[normalem]{ulem}
\renewcommand{\ULdepth}{1.8pt}
\renewcommand{\ULthickness}{0.2pt}
\contourlength{0.8pt}
\let\oldhref\href
\renewcommand{\href}[2]{\oldhref{#1}{\uline{#2}}}
\usepackage{fancyhdr}
\usepackage[english]{babel}
\usepackage{tabularx}
\input{glyphtounicode}


%----------FONT OPTIONS----------
% sans-serif
% \usepackage[sfdefault]{FiraSans}
% \usepackage[sfdefault]{roboto}
% \usepackage[sfdefault]{noto-sans}
% \usepackage[default]{sourcesanspro}

% serif
% \usepackage{CormorantGaramond}
% \usepackage{charter}


\pagestyle{fancy}
\fancyhf{} % clear all header and footer fields
\fancyfoot{}
\renewcommand{\headrulewidth}{0pt}
\renewcommand{\footrulewidth}{0pt}

% Adjust margins
\addtolength{\oddsidemargin}{-0.5in}
\addtolength{\evensidemargin}{-0.5in}
\addtolength{\textwidth}{1in}
\addtolength{\topmargin}{-.5in}
\addtolength{\textheight}{1.0in}

\urlstyle{same}

\raggedbottom
\raggedright
\setlength{\tabcolsep}{0in}

% Sections formatting
\titleformat{\section}{
    \vspace{-4pt}\scshape\raggedright\large
}{}{0em}{}[\color{black}\titlerule \vspace{-5pt}]

% Ensure that generate pdf is machine readable/ATS parsable
\pdfgentounicode=1

\pdfinfoomitdate=1
\pdftrailerid{}
\pdfsuppressptexinfo=-1

%-------------------------
% Custom commands
\newcommand{\resumeItem}[1]{
    \item\small{
            {#1 \vspace{-2pt}}
    }
}

\newcommand{\resumeSubheading}[4]{
    \vspace{-2pt}\item
    \begin{tabular*}{0.97\textwidth}[t]{l@{\extracolsep{\fill}}r}
    \textbf{#1} & #2 \\
    \textit{\small#3} & \textit{\small #4} \\
    \end{tabular*}\vspace{-7pt}
}

\newcommand{\resumeSubSubheading}[2]{
    \item
    \begin{tabular*}{0.97\textwidth}{l@{\extracolsep{\fill}}r}
    \textit{\small#1} & \textit{\small #2} \\
    \end{tabular*}\vspace{-7pt}
}

\newcommand{\resumeProjectHeading}[2]{
    \item
    \begin{tabular*}{0.97\textwidth}{l@{\extracolsep{\fill}}r}
    \small#1 & #2 \\
    \end{tabular*}\vspace{-7pt}
}

\newcommand{\resumeSubItem}[1]{\resumeItem{#1}\vspace{-4pt}}

\renewcommand\labelitemii{$\vcenter{\hbox{\tiny$\bullet$}}$}

\newcommand{\resumeSubHeadingListStart}{\begin{itemize}[leftmargin=0.15in, label={}]}
\newcommand{\resumeSubHeadingListEnd}{\end{itemize}}
\newcommand{\resumeItemListStart}{\begin{itemize}}
\newcommand{\resumeItemListEnd}{\end{itemize}\vspace{-5pt}}

%-------------------------------------------
%%%%%%  RESUME STARTS HERE  %%%%%%%%%%%%%%%%%%%%%%%%%%%%


\begin{document}

%----------HEADING----------
% \begin{tabular*}{\textwidth}{l@{\extracolsep{\fill}}r}
%   \textbf{\href{http://sourabhbajaj.com/}{\Large Sourabh Bajaj}} & Email : \href{mailto:sourabh@sourabhbajaj.com}{sourabh@sourabhbajaj.com}\\
%   \href{http://sourabhbajaj.com/}{http://www.sourabhbajaj.com} & Mobile : +1-123-456-7890 \\
% \end{tabular*}

    \begin{center}
        \textbf{\Huge \scshape Brandon Li} \\ \vspace{1pt}
        \small 978-245-5532 $|$
        \href{mailto:jobs@brandonli.me}{jobs@brandonli.me} $|$
        \href{https://brandonli.me}{brandonli.me} $|$
        \href{https://linkedin.com/in/brandonli28}{linkedin.com/in/brandonli28} $|$
        \href{https://github.com/PulseBeat02}{github.com/PulseBeat02}
    \end{center}


%-----------EDUCATION-----------
    \section{Education}
    \resumeSubHeadingListStart
    \resumeSubheading
    {University of California, Los Angeles (UCLA)}{Expected June {{GRAD_YEAR}}}
    {Bachelor of Science in Computer Science}{Los Angeles, CA}
    \resumeSubHeadingListEnd


%-----------EXPERIENCE-----------
    \section{Experience}
    \resumeSubHeadingListStart
    \resumeSubheading
    {Software Development Engineer Intern}{Sep 2026 -- Present}
    {Amazon (Amazon Web Services)}{Boston, MA}
    \resumeItemListStart
    \resumeItem{Extended Amazon’s Remote MCP server with 8 endpoints by integrating an internal work-management service (similar to Jira), enabling hundreds of thousands of employees to execute workflows directly from a chatbot.}
    \resumeItem{Built specialized AI agents for internal work-management workflows, allowing agent-to-agent (A2A) coordination for custom tasks such as summarization and board updates while handling thousands of production requests.}
    \resumeItem{Implemented A2A authentication and support for user-mocking agents for secure agent interactions and handoffs.}
    \resumeItem{Improved production reliability and debuggability by adding logging and over a hundred unit, wiring, and deployment tests, catching possible failures before release and validating end-to-end deployment behavior.}
    \resumeItemListEnd
    \vspace{1mm}
    \resumeSubheading
    {Software Engineering Intern}{June 2026 -- Sep 2026}
    {Google (YouTube)}{San Bruno, CA}
    \resumeItemListStart
    \resumeItem{Engineered 10-bit software AV1 HDR video playback for the YouTube Android app via low-level ARM/NEON SIMD vectorization in C++, improving the video quality and playback performance for millions of users.}
    \resumeItem{Defined per-device video resolution caps (e.g. 1080p) for dav1d software decoding, filtering higher-resolution formats out of the playback selection to prevent thermal throttling and battery drain on lower-end devices.}
    \resumeItem{Built a Java-based shader effect debug UI and OpenGL video pipeline to accelerate developer iteration speed.}
    \resumeItem{Integrated custom fragment shaders to allow video effects across thousands of DRM-protected shows and movies.}
    \resumeItem{Resolved memory leaks, migrated deprecated code, and improved format filtering across ExoPlayer and YouTube.}
    \resumeItemListEnd
    \vspace{1mm}
    \resumeSubheading
    {Software Engineering Intern}{June 2025 -- Sep 2025}
    {VideoLAN}{Remote}
    \resumeItemListStart
    \resumeItem{Built a shared audio dynamics core, cutting duplicate code by 20\% and binary size by 5\% to speed up downloads.}
    \resumeItem{Added ggml/OpenCV to VLC, ported 3 modules to C++, and upgraded to SAM3 to support text segmentation.}
    \resumeItem{Designed SAM2 object segmentation and YuNet face detection to be able to track over 100 concurrent targets.}
    \resumeItem{Implemented 5 different dithering algorithms to allow users to convert videos into GIF files at varying qualities.}
    \resumeItemListEnd
    \resumeSubHeadingListEnd


%-----------PROJECTS-----------
    \section{Projects}
    \resumeSubHeadingListStart
    \resumeProjectHeading
    {\href{https://github.com/PulseBeat02/yt-media-storage}{\textbf{yt-media-storage}} $|$ \emph{C++, FFmpeg, SIMD, OpenMP, Qt 6, libsodium, Wirehair}}{Feb 2026 -- Present}
    \resumeItemListStart
    \resumeItem{Built a C++ tool ({{YT_STORAGE_STARS}} stars, {{YT_STORAGE_FORKS}} forks) that can convert desktop files into YouTube videos for storage.}
    \resumeItem{Implemented SIMD intrinsics, OpenMP parallelization, and cache optimizations to speed up conversion by 100x.}
    \resumeItem{Designed a lightweight UI using Qt 6 and added libsodium encryption to enhance usability and security.}
    \resumeItem{Used CRC and Wirehair error correction to recover data corrupted by lossy video compression algorithms.}
    \resumeItem{Reached \#5 on Hacker News (Y Combinator) in 12 hours, drawing thousands of views and new contributors.}
    \resumeItem{Published a \href{https://www.youtube.com/watch?v=l03Os5uwWmk}{technical explanation video} on YouTube with over {{YT_IMPRESSIONS}} impressions, educating over {{YT_VIEWERS}} viewers.}
    \resumeItemListEnd
    \vspace{1mm}
    \resumeProjectHeading
    {\href{https://github.com/PulseBeat02/mcav}{\textbf{mcav}} $|$ \emph{Java, GLSL, libVLC, OpenCV, QEMU, VNC, TeamCity, Maven}}{June 2020 -- Present}
    \resumeItemListStart
    \resumeItem{Built a 12-module, 150K-line Java media framework ({{MCAV_STARS}} stars, {{MCAV_FORKS}} forks) whose unified player API streams real-time video from VLC, OpenCV, Chromium, VNC, and QEMU sources across four pluggable render backends.}
    \resumeItem{Designed a pure-GPU (GLSL fragment shader) video codec with O(1) per-pixel random access, improving BD-rate over 40\% with VMAF-scored optimization rounds to reach 10 ms/frame at 1080p for integrated graphics cards.}
    \resumeItem{Engineered a low-latency pipeline API with nearly 50 different documented video effects that can be applied.}
    \resumeItem{Self-hosted a TeamCity CI/CD pipeline that publishes releases to a Maven repository, gated with 2,500+ tests.}
    \resumeItemListEnd
    \resumeSubHeadingListEnd


%-----------PROGRAMMING SKILLS-----------
    \section{Technical Skills}
    \begin{itemize}[leftmargin=0.15in, label={}]
    \small{\item{
        \textbf{Languages}{: Java, C/C++, OCaml, JavaScript/TypeScript, Assembly} \\
        \textbf{Frameworks \& Libraries}{: Spring, FFmpeg, OpenCV, libVLC, ExoPlayer, Qt, OpenMP, React, Next.js, Tailwind} \\
        \textbf{Technologies}{: AWS, AI Agents, MCP, Git, CMake, Maven, TeamCity, OpenGL, Android, GLSL, AV1, NEON/SIMD} \\
        \textbf{Awards}{: Dean's Honors List}
    }}
    \end{itemize}


%-------------------------------------------
\end{document}
